% Déroulement du stage sur huit semaines. Chaque phase porte un livrable, et
% les jalons marquent les points où le travail devenait vérifiable.
\begin{ganttchart}[
    hgrid, vgrid,
    x unit=1.3cm,
    y unit title=0.62cm,
    y unit chart=0.58cm,
    title height=1,
    canvas/.style={draw=rule, line width=0.5pt, fill=none},
    bar top shift=0.16,
    group top shift=0.2,
    group height=0.36,
    group right shift=0,
    group left shift=0,
    milestone top shift=0.28,
    milestone height=0.44,
    bar incomplete/.append style={fill=mist},
    progress label text={},
  ]{1}{8}

\gantttitle{Phasage prévu \textcolor{slate!70}{— stage du 1\textsuperscript{er} juillet au 31 août 2026}}{8} \\
\gantttitlelist{1,...,8}{1} \\

% ------------------------------------------------------------------- corpus
\ganttgroup{Corpus}{1}{2} \\
\ganttbar{Extraction, découpage par article}{1}{1} \\
\ganttbar{Hiérarchie, métadonnées, JSONL}{2}{2} \\
\ganttmilestone{corpus reproductible}{2} \\

% ------------------------------------------------------------------- moteur
\ganttgroup{Moteur}{3}{5} \\
\ganttbar{Pipeline RAG bout-en-bout}{3}{3} \\
\ganttbar{Ancrage, citation, abstention}{4}{4} \\
\ganttbar{Recherche hybride, jeu de référence}{5}{5} \\
\ganttmilestone{réponses citées et sourcées}{4} \\

% --------------------------------------------------------- mesure et clôture
\ganttgroup{Mesure et clôture}{6}{8} \\
\ganttbar{Évaluation, analyse d'erreurs}{6}{6} \\
\ganttbar{Durcissement, reproductibilité}{7}{7} \\
\ganttbar{Rapport, démonstration, soutenance}{8}{8} \\
\ganttmilestone{banc de mesure rejouable}{6} \\

% Pas de \ganttlink : la numérotation « elem » compte aussi les groupes et les
% jalons, ce qui reliait des éléments arbitraires. L'enchaînement des phases se
% lit déjà dans leur succession hebdomadaire.

\end{ganttchart}
